% =====================================================================
%  CAPÍTULO 1 — Ejemplo: marco teórico (citas, notas al pie, definiciones)
% =====================================================================
\chapter{Título del capítulo 1}\label{cap:capitulo1}

\section{Introducción del capítulo}

% Citas automáticas desde referencias.bib:
%   \textcite{clave}  -> Autor (año)      \parencite{clave} -> (Autor, año)
La economía regional estudia la dimensión territorial de los procesos
económicos \parencite{capello2015}. \textcite{krugman1991} muestra que los
rendimientos crecientes pueden generar patrones de concentración
núcleo-periferia.\footnote{Las notas al pie se numeran automáticamente.}
\lipsum[5]

\section{Revisión de la literatura}

\subsection{Primer enfoque}

\lipsum[6]

\begin{definicion}\label{def:ejemplo}
Texto de la definición. Los términos en otro idioma se marcan con
\ingles{spillovers}.
\end{definicion}

\subsection{Segundo enfoque}

Para el caso mexicano, \textcite{esquivel1999} documenta convergencia regional,
en línea con \textcite{barro1992}. \lipsum[7]

\section{Consideraciones finales del capítulo}

\lipsum[8]
